\chapter*{尾声：把下一次修改留下来}
\addcontentsline{toc}{chapter}{尾声：把下一次修改留下来}

这本书最初想用一个漂亮故事收尾：一个学生写出有限元求解器，发现三个错误，最后说出一句关于“让桥站住”的话。故事完整、动人，也恰好证明全书的每一个观点。

可惜，在现有录音转写里找不到它。

如果一本讨论 AI 时代专业判断的书，需要靠一段无法核实的合成故事完成高潮，它已经背叛了自己的主题。因此，那个故事被删掉了。留下来的现场没有那么整齐，却更值得写。

2026 年 7 月 1 日的结题汇报持续近三个小时。学生展示了完成度差异很大的作品：有的模型能运行，却解释不清公式；有的游戏有新想法，但自动搭桥始终做不出来；有的小组发现 AI 把弯矩图画错，于是回到教材和课堂内容，逐个指定关键点；也有人坦率地说，继续迭代时把较好的半成品覆盖了。

这不是“AI 赋能教学”的整齐样板。它更像真实工程：工具有用，版本会坏，模型会错，时间有限，人必须作判断。

其中一个学生总结：

\begin{shiyu-inline}
``你就真的得是自己会才行，要不然你也不知道 AI 做得对不对。''

\hfill ——课程结题汇报，2026 年 7 月 1 日 00:25:20
\end{shiyu-inline}

这句话并不意味着学生应该先把所有知识学完，才有资格使用 AI。相反，项目把“还不会什么”暴露了出来。问题是课堂是否给学生第二次机会：回到来源，修改作品，解释变化，并把错误变成下一次判断的依据。

反馈研究把这一步称为关闭反馈回路：信息只有被学习者理解并用于后续行动，才成为反馈\cite{carless2019spirals}。这也解释了为什么本书反复强调修改记录。最终成品容易伪装，修改过程更接近学习发生的地方。

如果你是一位教师，读到这里不必立刻开发游戏、网站或 AI 代理。请先保存一份学生的初稿，再保存修改后的版本；问一句“你改了什么，依据是什么”；把回答写进下一次任务的评价标准。

如果你是一位学生，也请不要只保存最好看的结果。保留那张画错的图、那个无法运行的版本和你后来查到的依据。它们不是羞耻的残骸，而是专业判断长出来的年轮。

\enlargethispage{2\baselineskip}
桥梁工程里，验收不靠一张渲染图。学习也不该只靠一个最后答案。我们真正要留下的，是从第一次尝试到下一次修改之间那段可追溯的路。下一节课，就从那里开始。
